% Options for packages loaded elsewhere
\PassOptionsToPackage{unicode}{hyperref}
\PassOptionsToPackage{hyphens}{url}
\PassOptionsToPackage{dvipsnames,svgnames,x11names}{xcolor}
\documentclass[
  10pt,
]{article}
\usepackage{xcolor}
\usepackage[margin=0.82in]{geometry}
\usepackage{amsmath,amssymb}
\setcounter{secnumdepth}{-\maxdimen} % remove section numbering
\usepackage{iftex}
\ifPDFTeX
  \usepackage[T1]{fontenc}
  \usepackage[utf8]{inputenc}
  \usepackage{textcomp} % provide euro and other symbols
\else % if luatex or xetex
  \usepackage{unicode-math} % this also loads fontspec
  \defaultfontfeatures{Scale=MatchLowercase}
  \defaultfontfeatures[\rmfamily]{Ligatures=TeX,Scale=1}
\fi
\usepackage{lmodern}
\ifPDFTeX\else
  % xetex/luatex font selection
\fi
% Use upquote if available, for straight quotes in verbatim environments
\IfFileExists{upquote.sty}{\usepackage{upquote}}{}
\IfFileExists{microtype.sty}{% use microtype if available
  \usepackage[]{microtype}
  \UseMicrotypeSet[protrusion]{basicmath} % disable protrusion for tt fonts
}{}
\makeatletter
\@ifundefined{KOMAClassName}{% if non-KOMA class
  \IfFileExists{parskip.sty}{%
    \usepackage{parskip}
  }{% else
    \setlength{\parindent}{0pt}
    \setlength{\parskip}{6pt plus 2pt minus 1pt}}
}{% if KOMA class
  \KOMAoptions{parskip=half}}
\makeatother
\usepackage{longtable,booktabs,array}
\newcounter{none} % for unnumbered tables
\usepackage{calc} % for calculating minipage widths
% Correct order of tables after \paragraph or \subparagraph
\usepackage{etoolbox}
\makeatletter
\patchcmd\longtable{\par}{\if@noskipsec\mbox{}\fi\par}{}{}
\makeatother
% Allow footnotes in longtable head/foot
\IfFileExists{footnotehyper.sty}{\usepackage{footnotehyper}}{\usepackage{footnote}}
\makesavenoteenv{longtable}
\setlength{\emergencystretch}{3em} % prevent overfull lines
\providecommand{\tightlist}{%
  \setlength{\itemsep}{0pt}\setlength{\parskip}{0pt}}
\let\lanciusOriginalTableOfContents\tableofcontents
\renewcommand{\tableofcontents}{\begingroup\small\lanciusOriginalTableOfContents\endgroup}
\setlength{\emergencystretch}{2em}
\usepackage{bookmark}
\IfFileExists{xurl.sty}{\usepackage{xurl}}{} % add URL line breaks if available
\urlstyle{same}
\hypersetup{
  pdftitle={Lancius Mathematics Engine: Research and Development Plan},
  colorlinks=true,
  linkcolor={blue},
  filecolor={Maroon},
  citecolor={Blue},
  urlcolor={blue},
  pdfcreator={LaTeX via pandoc}}

\title{Lancius Mathematics Engine: Research and Development Plan}
\author{}
\date{2026-11-24}

\begin{document}
\maketitle

{
\setcounter{tocdepth}{1}
\tableofcontents
}
\clearpage

\textbf{Purpose:} establish a comprehensive, staged plan for making
Lancius's mathematical reasoning more accurate, explainable, and
resource-efficient.

\textbf{Status:} research plan. It records directions to investigate and
criteria for choosing them; it does not claim that the proposed
algorithms have been implemented or benchmarked.

\textbf{Starting-point snapshot:} 2026-11-24. The I6 baseline described
below was recorded on that date. V1.2-RC3.4 publishes a later model with a
different sparse input schema and held-out evaluation contract; see
\href{../generation_1/I6/README.md}{\texttt{generation\_1/I6/README.md}}
for its implementation and measurements.

\section{1. Goal}\label{goal}

Build a mathematics engine that can solve and check supported problems
with exact arithmetic, explain each transformation, and reject or
abstain on work outside its supported domain. Keep the trusted checker
small. Let more complicated routines propose answers or proof steps,
then require the checker to validate a compact certificate.

For known, formalized mathematics, use the domain algorithm as the
source of the result. Use learned or heuristic methods for parsing,
candidate generation, and prioritizing work. Keep them outside the final
trust boundary.

The aim is a strong collection of dependable methods for specific
mathematical domains. A single procedure that decides every mathematical
theorem is not an achievable target: first-order logic is undecidable in
general
(\href{https://plato.stanford.edu/entries/goedel-incompleteness/}{Church's
theorem}). Each solver therefore needs a stated domain, assumptions,
resource limit, and failure mode.

\section{2. Current starting point}\label{current-starting-point}

\begin{itemize}
\tightlist
\item
  \href{../tools/prm/algos.py}{\texttt{tools/prm/algos.py}} defines the
  original mathematical workload with exact \texttt{Fraction} values and
  canonical derivation steps. The verifier lineage has expanded this
  into 150 algorithms across algebra, calculus, linear algebra, and
  number theory.
\item
  \href{../tools/prm/verifier5.py}{\texttt{tools/prm/verifier5.py}}
  computes exact step agreement and structural agreement, then emits
  measurements used by the I6 training corpus.
\item
  The 2026-11-24 I6 baseline used a 65 → 4096 → 2560 → 5 network with
  10,771,461 parameters per family expert. RC3.4 uses 65 → 128 → 64 → 5
  with seven semantic inputs and 17,029 parameters per expert, fresh
  initialization, and algorithm-held-out splits. See
  \href{../generation_1/I6/README.md}{\texttt{generation\_1/I6/README.md}}.
\item
  The I6 feature schema changed coordinate 57 from \texttt{mean\_rel} to
  \texttt{n\_wrong}; the earlier I1--I3 models retain the old meaning.
  RC3.4 has a separate sparse feature-generation path
  and rejects incompatible I3 warm starts. Preserve this separation in
  any new corpus, checkpoint, or evaluator.
\item
  The retained I3 normalizer is not a fitted I6 z-score for
  \texttt{n\_wrong}. Existing model artifacts depend on their training
  transform. Future model formats should identify the feature schema and
  normalizer explicitly.
\item
  Prior saved-model continuous-error claims were audited and withdrawn.
  Any new algorithm should be compared with exact, per-row measurements
  on a fixed evaluation set; no score should be inferred from argmax
  accuracy or from \texttt{MAE²}.
\end{itemize}

\section{3. Target architecture}\label{target-architecture}

\begin{center}
\fbox{\begin{minipage}{0.92\linewidth}
\centering
\textbf{Problem or candidate derivation}\\[0.35em]
$\downarrow$\\[0.25em]
\textbf{Typed expression tree with domain and assumptions}\\[0.35em]
$\downarrow$\\[0.25em]
\textbf{Bounded exact solver}\\[-0.1em]
{\small (heuristic or learned proposer may supply candidate steps)}\\[0.35em]
$\downarrow$\\[0.25em]
\textbf{Result and proof certificate} $\longrightarrow$ \textbf{Independent checker}\\[0.35em]
$\downarrow$\\[0.25em]
\textbf{Verified result or explicit unknown}
\end{minipage}}
\end{center}

\subsection{Trust boundary}\label{trust-boundary}

\begin{enumerate}
\def\labelenumi{\arabic{enumi}.}
\tightlist
\item
  \textbf{Parser and type checker:} determine whether a symbol denotes
  an integer, rational, real approximation, matrix, modular residue, or
  proposition. Preserve the distinction between these domains.
\item
  \textbf{Solver:} simplify, solve, or construct a proof. It may be
  optimized, heuristic, or eventually external.
\item
  \textbf{Certificate:} state the theorem, assumptions, rule
  identifiers, substitutions, premises, and conclusion for each step.
\item
  \textbf{Checker:} independently confirm each local inference and the
  final result using exact operations. A timeout, unsupported construct,
  failed precondition, or invalid certificate returns \texttt{unknown}
  or \texttt{invalid}; it never silently becomes a guessed answer.
\end{enumerate}

Lean uses a related arrangement: tactics construct proof terms and a
small kernel checks them. This is a useful architecture to study,
without requiring Lancius to adopt Lean as a runtime dependency.
\href{https://lean-lang.org/doc/reference/latest/Elaboration-and-Compilation/}{Lean
kernel and proof-checking reference}

\section{4. Foundational mathematics to
formalize}\label{foundational-mathematics-to-formalize}

The first deliverable should be a versioned \textbf{mathematical
contract}: exact definitions, preconditions, accepted input forms, and
expected outputs for every supported rule. These are the rules the
algorithms must preserve.

\subsection{4.1 Logic and equality}\label{logic-and-equality}

Implement the inference rules used by every later domain:

\begin{itemize}
\tightlist
\item
  Equality: reflexivity, symmetry, transitivity, and congruence.
\item
  Substitution: replace equals with equals only in a well-typed context.
\item
  Propositional logic: conjunction/disjunction introduction and
  elimination, implication, modus ponens, negation, and contradiction.
\item
  Quantifiers: variable binding, capture-avoiding substitution, and
  explicit universal instantiation where supported.
\item
  Equivalence: record both directions when a transformation is
  reversible; record a one-way implication when it is not.
\end{itemize}

Every rule has a small checker and a certificate form. A proof step
should identify the rule used rather than merely show two strings that
appear equal.

\subsection{4.2 Arithmetic domains}\label{arithmetic-domains}

Represent the domain in the expression itself:

\begin{itemize}
\tightlist
\item
  Integers \(\mathbb Z\).
\item
  Rationals \(\mathbb Q\), represented as normalized \(p/q\), with
  \(q>0\).
\item
  Modular integers \(\mathbb Z/n\mathbb Z\), with an explicit modulus.
\item
  Real numbers represented by exact algebraic objects or by certified
  intervals when needed.
\item
  Floating point as a separate finite-precision domain, never as an
  implicit substitute for exact real arithmetic.
\end{itemize}

Core laws include the semiring/ring laws, field laws when division is
allowed, order laws, and the domain-specific rules for modular
arithmetic. Division, square roots, logarithms, and inequality
transformations must carry their preconditions.

Examples of preconditions that must be represented:

\begin{itemize}
\tightlist
\item
  Dividing both sides by \(a\) requires \(a\ne0\).
\item
  Multiplying an inequality by a negative value reverses its direction.
\item
  Cancelling a factor requires proving it is nonzero.
\item
  \(\sqrt{x}\) over the reals requires \(x\ge0\).
\item
  Reducing a congruence or taking a modular inverse depends on the
  modulus and, for inverses, coprimality.
\end{itemize}

\subsection{4.3 Algebra identities and
equations}\label{algebra-identities-and-equations}

Prioritize identities already represented in the corpus, then expand by
demand:

\begin{itemize}
\tightlist
\item
  Linear equation: \(ax+b=c\), with explicit cases for \(a=0\).
\item
  Quadratic formula and discriminant: \(D=b^2-4ac\),
  \(x=(-b\pm\sqrt D)/(2a)\), with \(a\ne0\) and real-domain checks.
\item
  Difference of squares, perfect-square trinomials, sums/differences of
  cubes, binomial theorem, and polynomial expansion.
\item
  Arithmetic and geometric sequences and finite sums.
\item
  Polynomial equality and factorization over a declared coefficient
  domain.
\item
  Rational-expression equality, with denominator restrictions retained.
\end{itemize}

\subsection{4.4 Number theory}\label{number-theory}

\begin{itemize}
\tightlist
\item
  Euclidean algorithm for \(\gcd(a,b)\).
\item
  Extended Euclidean algorithm for Bézout coefficients
  \(ax+by=\gcd(a,b)\).
\item
  Divisibility, prime-factor checks within a bounded range, congruences,
  and modular inverses.
\item
  Modular exponentiation by repeated squaring.
\item
  Integer square-root and perfect-square checks.
\item
  Later: stronger primality/factorization algorithms only if a real
  workload requires them.
\end{itemize}

\subsection{4.5 Linear algebra}\label{linear-algebra}

\begin{itemize}
\tightlist
\item
  Row operations and their effect on the solution set.
\item
  Gaussian elimination for small rational systems, including pivoting,
  rank, consistency, and free variables.
\item
  Determinant, inverse, nullspace, dot product, matrix-vector product,
  and matrix identities.
\item
  For integer matrices, evaluate fraction-free Bareiss elimination
  before standard rational elimination. Its updates use exact divisions
  and can avoid large intermediate fractions; FLINT's algorithm notes
  give a useful implementation reference, alongside Bareiss's original
  paper.
  \href{https://doi.org/10.1090/S0025-5718-1968-0226829-0}{Bareiss
  paper},
  \href{https://flintlib.org/doc/algorithms/linear-algebra/fraction-free.html}{FLINT
  fraction-free elimination}.
\end{itemize}

\subsection{4.6 Calculus}\label{calculus}

Start with the parts having small exact checkers:

\begin{itemize}
\tightlist
\item
  Derivative of constants, sums, products, quotients, integer powers,
  and compositions.
\item
  Polynomial derivatives and exact polynomial antiderivatives.
\item
  Definite polynomial integrals by evaluating an antiderivative at the
  endpoints.
\item
  Basic limit rules only with explicit side conditions.
\item
  Fundamental theorem of calculus in a formal statement with its
  continuity/integrability assumptions.
\end{itemize}

General antiderivatives, transcendental identities, and unrestricted
limits are a later research area. They require stronger domain and
branch-cut semantics than the polynomial cases.

\subsection{4.7 Probability and
inequalities}\label{probability-and-inequalities}

\begin{itemize}
\tightlist
\item
  Counting and finite sums using exact integers/rationals.
\item
  Conditional probability, complement, total probability, Bayes'
  theorem, and expectation for finite distributions.
\item
  Order transitivity, triangle inequality, sign-aware multiplication,
  AM--GM, Cauchy--Schwarz, and Pythagorean identities.
\item
  Separate a proof of an inequality from a numerical approximation to
  its two sides.
\end{itemize}

\section{5. Candidate algorithms and when to use
them}\label{candidate-algorithms-and-when-to-use-them}

{\def\LTcaptype{none} % do not increment counter
\begin{longtable}[]{@{}
  >{\raggedright\arraybackslash}p{(\linewidth - 6\tabcolsep) * \real{0.2500}}
  >{\raggedright\arraybackslash}p{(\linewidth - 6\tabcolsep) * \real{0.2500}}
  >{\raggedright\arraybackslash}p{(\linewidth - 6\tabcolsep) * \real{0.2500}}
  >{\raggedright\arraybackslash}p{(\linewidth - 6\tabcolsep) * \real{0.2500}}@{}}
\toprule\noalign{}
\begin{minipage}[b]{\linewidth}\raggedright
Method
\end{minipage} & \begin{minipage}[b]{\linewidth}\raggedright
Best use
\end{minipage} & \begin{minipage}[b]{\linewidth}\raggedright
Accuracy and cost considerations
\end{minipage} & \begin{minipage}[b]{\linewidth}\raggedright
Priority
\end{minipage} \\
\midrule\noalign{}
\endhead
\bottomrule\noalign{}
\endlastfoot
Exact rational evaluation & Arithmetic and finite derivations & Exact
results; cost grows with integer bit-length. & Immediate \\
Canonical polynomial form & Polynomial identities and equality & Store
monomials in a fixed variable order and combine coefficients. Expression
expansion can grow rapidly, so cap terms and bit-length. & Immediate \\
Oriented rewrite rules & Simplifying expressions and checking steps &
Cheap and explainable when each rewrite decreases a well-founded
measure. Unrestricted bidirectional rules can loop. & Immediate \\
Horner evaluation & Evaluating polynomials & Linear in polynomial degree
with fewer multiplications and less intermediate growth than naive
evaluation. & Immediate \\
Euclid / extended Euclid & GCD, divisibility, modular inverses & Very
small memory; logarithmic number of division steps in input magnitude. &
Immediate \\
Repeated squaring & Powers and modular powers & \(O(\log n)\)
multiplications for exponent \(n\); check overflow or work in an
explicit modular domain. & Immediate \\
Linear/quadratic solvers & Common equation classes & Constant-size
logic; enumerate degenerate coefficients and exact-root cases. &
Immediate \\
Sparse exact Gaussian elimination & Small systems over \(\mathbb Q\) &
Pivoting and sparsity matter; rational numerators/denominators may grow.
& High \\
Bareiss fraction-free elimination & Integer matrices & Exact divisions
can control fraction growth; requires correct pivot bookkeeping. &
High \\
Interval arithmetic & Certified bounds for irrational or floating
quantities & A result is an enclosure, not an exact equality. Use
directed rounding and report interval width. & High when needed \\
SMT / DPLL(T) & Conjunctions of logic, equality, and supported
arithmetic constraints & Reuses specialized decision procedures. Bound
solver time and validate proof/counterexample output. SMT-LIB defines
logics including linear integer/real arithmetic and bit vectors.
\href{https://smt-lib.org/language.shtml}{SMT-LIB standard},
\href{https://smt-lib.org/logics.shtml}{supported logics}. & Medium \\
Equality saturation / e-graphs & Searching among equivalent expression
forms & Can find useful simplifications without destructive rewrite
order; e-graphs can grow quickly. Apply node/time limits and extract a
proof path. \href{https://doi.org/10.1145/3434304}{egg paper}. &
Later \\
Polynomial GCD / Gröbner bases & Multivariate polynomial equations and
ideal membership & Potentially expensive; do not put in the small
trusted core. Use only for demonstrated workloads, with checkable
outputs. & Later \\
\end{longtable}
}

For matrix and polynomial code, compare custom exact routines against
mature arithmetic libraries such as GMP or FLINT before deciding whether
to depend on them. The comparison must include build footprint,
portability, license, memory, and certificate/checker requirements---not
just speed.

\subsection{Questions to answer before choosing an
implementation}\label{questions-to-answer-before-choosing-an-implementation}

\begin{enumerate}
\def\labelenumi{\arabic{enumi}.}
\tightlist
\item
  \textbf{Is Lancius checking a known procedure or judging an unknown
  derivation?} For a known procedure, compare direct execution plus a
  proof certificate against the learned verifier. For unknown
  procedures, keep theorem search separate from certificate checking.
\item
  \textbf{Which mathematical domain is intended for each claim?}
  Establish whether each value is in the integers, rationals, reals,
  modular arithmetic, or floating point before selecting
  transformations.
\item
  \textbf{Which current failures are arithmetic failures, and which are
  parsing, semantic, feature, or training failures?} Preserve these as
  distinct failure codes and benchmark each separately.
\item
  \textbf{When does exact arithmetic become too expensive?} Measure
  operand bits and expression growth; choose a declared cutoff and a
  certified interval fallback rather than silent floating conversion.
\item
  \textbf{Which normal form works for each family?} Compare sparse
  polynomial collection, rational-function normalization, ordered
  rewriting, and later e-graph search. Do not force one normalizer on
  every domain.
\item
  \textbf{Which linear constraint fragment covers the next workload?}
  Compare simple substitution/elimination with Fourier--Motzkin or a
  bounded SMT solver. General integer/nonlinear arithmetic has different
  completeness and cost properties from rational linear arithmetic.
\item
  \textbf{Which library dependencies are acceptable?} Compare custom
  integers with GMP/FLINT for C integration; compare solver libraries
  only after certificate export and independent checking are
  demonstrated.
\item
  \textbf{What is the resource budget?} Set a per-problem time, memory,
  integer-bit, AST-node, and certificate-size limit before enabling more
  advanced search.
\end{enumerate}

\subsection{Prototype and integration
policy}\label{prototype-and-integration-policy}

\begin{itemize}
\tightlist
\item
  Prototype mathematical semantics in \texttt{tools/prm} using the
  existing exact-rational workflow and explicit proof records.
\item
  Keep each algorithm callable independently so it can be cross-checked
  without training or running the full I6 pipeline.
\item
  Port stable routines to C only when the runtime needs them and the
  certificate contract is settled.
\item
  Preserve a small C checker for certificates even if a larger solver is
  added in Python, Lean, an SMT solver, or a computer-algebra library.
\item
  Add a model or heuristic only for work that cannot be decided cheaply
  by a supported exact routine, such as candidate parsing or tactic
  selection.
\item
  Version every rule set and certificate format. Version feature
  meanings and normalizers separately for any learned fallback.
\end{itemize}

\section{6. Efficiency work to
investigate}\label{efficiency-work-to-investigate}

\subsection{6.1 Exact rational
arithmetic}\label{exact-rational-arithmetic}

\begin{itemize}
\tightlist
\item
  Benchmark eager reduction after every operation against
  delayed/cross-cancelled reduction.
\item
  For multiplication \((a/b)(c/d)\), cancel \(\gcd(|a|,d)\) and
  \(\gcd(|c|,b)\) before multiplying.
\item
  For addition, use the denominator gcd to form a
  least-common-denominator path instead of multiplying both denominators
  blindly.
\item
  Track numerator/denominator bit-length, gcd count, and total
  big-integer allocations.
\item
  If fixed-width C integers are used, detect overflow before it occurs.
  Never wrap and continue as if the result were exact.
\item
  Compare a bounded built-in integer representation, GMP, and FLINT on
  the actual Lancius corpus before choosing a backend.
\end{itemize}

\subsection{6.2 Expression
representation}\label{expression-representation}

\begin{itemize}
\tightlist
\item
  Use a typed abstract syntax tree (AST) with explicit literals,
  symbols, operations, relations, and assumptions.
\item
  Hash-cons repeated subexpressions into a DAG; memoize exact
  evaluations and normal forms.
\item
  Represent polynomials sparsely as
  \texttt{(monomial,\ rational\ coefficient)} pairs with canonical
  ordering.
\item
  Normalize signs, rational constants, and associative/commutative
  operators in one defined pass.
\item
  Preserve original source spans and rewrite provenance so every
  normalized result can be explained.
\item
  Add hard limits for AST nodes, polynomial terms, proof length, integer
  bits, and wall-clock work.
\end{itemize}

\subsection{6.3 Rewrite safety}\label{rewrite-safety}

\begin{itemize}
\tightlist
\item
  Give each rewrite a formal rule name, typed signature, assumptions,
  and proof obligation.
\item
  Orient basic rewrites toward a chosen normal form; prove or
  structurally enforce termination for each rule set.
\item
  Keep transformations that are only implications distinct from
  equivalences.
\item
  Never simplify across a denominator, square root, logarithm, absolute
  value, or inequality without retaining the relevant condition.
\item
  Investigate e-graphs only after deterministic normal forms and rule
  certificates work. Equality saturation is a search structure, not a
  proof checker by itself.
\end{itemize}

\subsection{6.4 Linear algebra and polynomial
growth}\label{linear-algebra-and-polynomial-growth}

\begin{itemize}
\tightlist
\item
  Compare rational elimination, fraction-free Bareiss, and modular
  methods on realistic matrix sizes and coefficient lengths.
\item
  Select pivots that limit coefficient growth; preserve row permutation
  and sign in certificates.
\item
  Use sparse storage when matrices are sparse; do not pay dense costs
  for zero entries.
\item
  For polynomial expressions, measure term explosion before considering
  Gröbner bases or general expansion.
\item
  Consider modular reconstruction for large integer determinants only
  after small-matrix Bareiss performance is measured.
\end{itemize}

\subsection{6.5 Certified approximate
mathematics}\label{certified-approximate-mathematics}

\begin{itemize}
\tightlist
\item
  Use exact rationals whenever the problem is rational.
\item
  Use outward-rounded intervals for certified bounds, not a floating
  tolerance pretending to prove equality.
\item
  Return ``proved by interval'' only when the enclosure establishes the
  required sign or containment; return ``inconclusive'' when it
  straddles the boundary.
\item
  Treat algebraic roots as a later type: polynomial plus isolating
  interval, with exact root-count evidence if required.
\end{itemize}

\section{7. Work phases and decision
gates}\label{work-phases-and-decision-gates}

\subsection{Phase 0 --- Freeze the mathematical
contract}\label{phase-0-freeze-the-mathematical-contract}

\textbf{Deliverables:} theorem/rule catalogue; type/domain definitions;
versioned feature schema; certificate format; unsupported-case behavior.

\textbf{Gate:} every current corpus label and every mathematical
operation has a precise definition and explicit preconditions. I1--I3
and I6 feature schemas, normalizers, and saved-model metadata remain
distinguishable.

\subsection{Phase 1 --- Exact arithmetic
kernel}\label{phase-1-exact-arithmetic-kernel}

Implement normalized rationals, exact comparison, integer division, gcd,
powers, perfect-square checks, and explicit resource limits. Record
bit-growth and allocation counters.

\textbf{Gate:} all arithmetic operations either return an exact result
or an explicit failure. No overflow, failed division, or malformed
rational can produce an accepted answer.

\subsection{Phase 2 --- Proof certificates and small
checker}\label{phase-2-proof-certificates-and-small-checker}

Implement certificate records for arithmetic, equality, substitution,
linear solving, and existing canonical steps. Keep the checker
independent of the routine that generated the answer.

\textbf{Gate:} deliberately corrupted certificates are rejected. Valid
certificates round-trip through serialization and are accepted by the
checker.

\subsection{Phase 3 --- High-value equation and identity
modules}\label{phase-3-high-value-equation-and-identity-modules}

Implement linear equations, quadratics, polynomial identities, binomial
expansion, arithmetic/geometric sums, and rational-expression checks.
Compare canonical forms with step-by-step certificates.

\textbf{Gate:} exact results for all supported cases; explicit unknown
for unsupported roots, invalid domains, and resource-limit cases.

\subsection{Phase 4 --- Number theory and small linear
algebra}\label{phase-4-number-theory-and-small-linear-algebra}

Implement Euclid, extended Euclid, modular exponentiation/inverse,
small-prime/perfect-square checks, row reduction, rank, determinant, and
small-system solving. Benchmark Bareiss against ordinary rational
elimination.

\textbf{Gate:} independently cross-check each solver; certificates
preserve divisibility, pivot, rank, and permutation facts.

\subsection{Phase 5 --- Calculus, probability, and
inequalities}\label{phase-5-calculus-probability-and-inequalities}

Implement exact finite probability first, then derivative rules and
polynomial integration. Add inequality transformations with sign/domain
conditions. Add certified intervals only where exact forms are
insufficient.

\textbf{Gate:} each theorem checker validates its assumptions and exact
result; interval cases report bounds and inconclusive outcomes honestly.

\subsection{Phase 6 --- Bounded
automation}\label{phase-6-bounded-automation}

Prototype SMT for supported linear arithmetic and equality constraints.
Prototype equality saturation for expression simplification. Measure
proof/certificate extraction costs and impose hard resource limits.

\textbf{Gate:} every accepted automated result is checked by the Lancius
checker or a separately trusted checker. A timeout is not a proof of
falsehood.

\subsection{Phase 7 --- Optional model
assistance}\label{phase-7-optional-model-assistance}

Keep any learned model as a parser, tactic selector, or candidate-step
proposer. Feed its suggestions into the exact checker. Retain an
explicit fallback to exact routines and an abstention result.

\textbf{Gate:} learned components cannot directly issue a mathematically
verified verdict without a valid certificate.

\section{8. Verification and benchmark
protocol}\label{verification-and-benchmark-protocol}

This section specifies future evaluation; it does not report test runs.

\subsection{8.1 Correctness tests}\label{correctness-tests}

\begin{itemize}
\tightlist
\item
  Exhaustively enumerate small integer/rational inputs for each bounded
  rule.
\item
  Generate random exact inputs independently of the solver
  implementation.
\item
  Mutate one proof step at a time: sign, operator, denominator, pivot,
  exponent, omitted step, or assumption.
\item
  Include degeneracies: zero coefficients, zero denominators, negative
  values, repeated roots, singular matrices, empty sums, modulus 1,
  non-coprime modular inverses, missing values, and long derivations.
\item
  Use metamorphic checks: commutative rearrangement, equivalent fraction
  scaling, applying an identity then its inverse, and row operations
  that preserve a linear system.
\item
  Cross-check selected cases with an independently implemented exact
  system. Never use the same generator logic as both the sole oracle and
  the implementation under test.
\end{itemize}

\subsection{8.2 Dataset and generalization
checks}\label{dataset-and-generalization-checks}

\begin{itemize}
\tightlist
\item
  Preserve a fixed clean test set and exact row-level metrics.
\item
  Report separate results for known algorithms, held-out parameter
  ranges, held-out algorithms, and held-out families.
\item
  Distinguish step correctness, final-answer correctness, structural
  correctness, proof validity, and unsupported-case abstention.
\item
  Evaluate the exact engine and I6 on identical cases with each model
  receiving its own versioned feature schema and normalization.
\item
  Avoid using the synthetic distribution as evidence for unrestricted
  human-written proof competence.
\end{itemize}

\subsection{8.3 Efficiency measurements}\label{efficiency-measurements}

Record, by algorithm family and input size:

\begin{itemize}
\tightlist
\item
  Median and 95th-percentile runtime.
\item
  Peak memory and allocation count.
\item
  Integer/rational operand bit-length and polynomial term count.
\item
  Number of rewrite steps, solver decisions, and certificate nodes.
\item
  Checker time separately from solver time.
\item
  Cold and warm cache results.
\item
  Model parameters and inference operations for any learned fallback.
\end{itemize}

\subsection{8.4 Acceptance rules}\label{acceptance-rules}

\begin{itemize}
\tightlist
\item
  \textbf{Soundness:} no false acceptance in the independent bounded
  test suite; all supported outputs carry checkable evidence.
\item
  \textbf{Coverage:} unsupported inputs return a clear
  unknown/unsupported result, with no forced class.
\item
  \textbf{Accuracy:} report exact answer and proof-step metrics
  separately; no MAE/MSE inferred from class accuracy.
\item
  \textbf{Efficiency:} retain a method only if it improves latency,
  memory, coverage, or clarity for a measured workload.
\item
  \textbf{Reproducibility:} version theorem rules, domains, feature
  schemas, normalizers, corpus generation, and benchmarks.
\end{itemize}

\section{9. Risks and boundaries}\label{risks-and-boundaries}

\begin{itemize}
\tightlist
\item
  \textbf{Undecidability:} no general procedure can settle all
  first-order mathematical claims. Specialize to decidable fragments and
  bounded algorithms.
\item
  \textbf{Domain confusion:} an integer, rational, real approximation,
  and modular residue obey different laws. Implicit conversion can make
  a valid-looking but false proof.
\item
  \textbf{Unsound cancellation:} algebraic cancellation and inequality
  transformations require nonzero/sign assumptions.
\item
  \textbf{Expression swell:} exactness can consume time and memory when
  numerators, denominators, or polynomial term counts grow. Use limits
  and explicit unknown results.
\item
  \textbf{Shared-bug validation:} agreement with the same generator or
  checker is weak evidence. Keep independent oracles and mutation tests.
\item
  \textbf{Feature-schema drift:} dimensions alone do not identify
  feature meaning. Store a schema ID and normalizer with every trained
  artifact.
\item
  \textbf{Overclaiming:} success on the generated 150-algorithm corpus
  establishes competence on that formalized workload, not general
  mathematical reasoning.
\end{itemize}

\section{10. Research references}\label{research-references}

These are starting references for the proposed workstreams, not
endorsements of adopting the referenced software as a dependency.

\begin{enumerate}
\def\labelenumi{\arabic{enumi}.}
\tightlist
\item
  \href{https://lean-lang.org/doc/reference/latest/Elaboration-and-Compilation/}{Lean
  language reference: kernel and proof checking} --- small trusted
  checker architecture.
\item
  \href{https://smt-lib.org/language.shtml}{SMT-LIB language standard}
  and \href{https://smt-lib.org/logics.shtml}{logic catalogue} ---
  standard encodings and supported arithmetic/bit-vector fragments.
\item
  Erwin H. Bareiss,
  \href{https://doi.org/10.1090/S0025-5718-1968-0226829-0}{``Sylvester's
  Identity and Multistep Integer-Preserving Gaussian Elimination''},
  1968 --- fraction-free elimination.
\item
  \href{https://flintlib.org/doc/algorithms/linear-algebra/fraction-free.html}{FLINT
  fraction-free elimination notes} --- practical pivoting, exact
  division, and integer-growth details; cross-check implementation
  claims against the original paper.
\item
  Max Willsey et al., \href{https://doi.org/10.1145/3434304}{``egg: Fast
  and Extensible Equality Saturation''}, POPL 2021 --- equality
  saturation and e-graphs for rewrite search.
\item
  \href{https://plato.stanford.edu/entries/goedel-incompleteness/}{Church's
  theorem overview} --- boundary on general theorem decision procedures.
\end{enumerate}

\section{11. Immediate next actions}\label{immediate-next-actions}

\begin{enumerate}
\def\labelenumi{\arabic{enumi}.}
\tightlist
\item
  Freeze and document the mathematical domains and rule preconditions
  for the existing corpus.
\item
  Define a certificate format and a minimal independent checker for
  exact rational steps.
\item
  Implement and benchmark cross-cancelled rational arithmetic and sparse
  polynomial normalization.
\item
  Reproduce the current 150-algorithm workload with exact rule-based
  outputs and compare against I6.
\item
  Add Euclid/extended Euclid, modular exponentiation, linear solving,
  and small-matrix Bareiss as the first new algorithm set.
\item
  Build the independent correctness and resource benchmark suite before
  pursuing SMT, e-graphs, Gröbner bases, or broader calculus.
\end{enumerate}

\section{Appendix A --- Individual theorem and equation
registry}\label{appendix-a-individual-theorem-and-equation-registry}

This registry turns the broad math families into concrete rules with
domains and checks. These are mathematical targets for a versioned
engine contract. A formula appearing in the current synthetic corpus may
be tested only on a small parameter range or a single instance; Appendix
B lists those source descriptions individually.

\subsection{A.1 First rules for the trusted exact
kernel}\label{a.1-first-rules-for-the-trusted-exact-kernel}

These rules are the highest priority because every later solver depends
on them.

\begin{itemize}
\tightlist
\item
  \textbf{Equality reflexivity:} \(a=a\). Check that both sides have the
  same typed expression.
\item
  \textbf{Equality symmetry:} \(a=b\Rightarrow b=a\). Reference the
  earlier equality and reverse its sides.
\item
  \textbf{Equality transitivity:} \(a=b\land b=c\Rightarrow a=c\). The
  middle terms must be the same typed expression.
\item
  \textbf{Congruence:} \(a=b\Rightarrow f(a)=f(b)\), only where \(f\) is
  defined for both values in the declared domain.
\item
  \textbf{Rational normalization:} represent \(p/q\) with \(q>0\) and
  \(\gcd(|p|,q)=1\); normalize zero to \(0/1\). Reject \(q=0\).
\item
  \textbf{Rational equality:} for \(b,d\ne0\), \(a/b=c/d\iff ad=bc\).
  Check by integer cross multiplication.
\item
  \textbf{Rational addition and multiplication:}
  \(a/b+c/d=(ad+bc)/(bd)\) and \((a/b)(c/d)=ac/(bd)\), with \(b,d\ne0\).
  Reduce without changing the value.
\item
  \textbf{Rational division:} \((a/b)/(c/d)=ad/(bc)\), requiring
  \(b,d,c\ne0\). The divisor must be nonzero.
\item
  \textbf{Euclidean division:} for integer \(a\) and positive integer
  \(b\), there are unique integers \(q,r\) with \(a=qb+r\) and
  \(0\le r<b\).
\item
  \textbf{GCD remainder step:} \(\gcd(a,b)=\gcd(b,a\bmod b)\) for
  \(b\ne0\); stop at remainder zero and return a nonnegative gcd.
\item
  \textbf{Bézout identity:} extended Euclid returns \(x,y\) with
  \(ax+by=\gcd(a,b)\). Verify this equality directly.
\item
  \textbf{Divisibility cancellation:} if \(d\mid ab\) and
  \(\gcd(d,a)=1\), then \(d\mid b\). Never cancel without the
  coprimality condition.
\item
  \textbf{Congruence definition:} \(a\equiv b\pmod m\iff m\mid(a-b)\),
  for integer \(m>0\).
\item
  \textbf{Modular inverse existence:} \(a^{-1}\pmod m\) exists exactly
  when \(\gcd(a,m)=1\). Verify a claimed inverse with
  \(ax\equiv1\pmod m\).
\item
  \textbf{Order under multiplication:} if \(x<y\), multiplying by
  \(c>0\) preserves the inequality; multiplying by \(c<0\) reverses it;
  multiplying by zero makes both sides equal.
\end{itemize}

\subsection{A.2 Algebra equations and
identities}\label{a.2-algebra-equations-and-identities}

\begin{itemize}
\tightlist
\item
  \textbf{Linear equation (\texttt{linear\_solve}):} for \(ax+b=c\), if
  \(a\ne0\), the unique solution is \(x=(c-b)/a\). If \(a=0\), there are
  all solutions when \(b=c\) and no solutions otherwise; the current
  routine covers only \(a\ne0\).
\item
  \textbf{Two-equation sum/difference system
  (\texttt{substitution\_system}):} \(x+y=S,\ x-y=D\) gives
  \(x=(S+D)/2,\ y=(S-D)/2\). Substitute both answers into both original
  equations.
\item
  \textbf{Quadratic formula (\texttt{quadratic\_formula}):} for
  \(ax^2+bx+c=0\) and \(a\ne0\), let \(\Delta=b^2-4ac\); then
  \(x=(-b\pm\sqrt\Delta)/(2a)\). Over the reals, \(\Delta<0\) means no
  real roots, \(\Delta=0\) one repeated root, and \(\Delta>0\) two
  roots. Current generated cases use rational perfect-square
  discriminants.
\item
  \textbf{Quadratic Vieta identities:} if roots are \(r_1,r_2\), then
  \(r_1+r_2=-b/a\) and \(r_1r_2=c/a\), for \(a\ne0\).
\item
  \textbf{Cubic Vieta identities (\texttt{vieta\_cubic}):} if
  \(a(x-r_1)(x-r_2)(x-r_3)=ax^3+bx^2+cx+d\), then \(r_1+r_2+r_3=-b/a\),
  \(r_1r_2+r_1r_3+r_2r_3=c/a\), and \(r_1r_2r_3=-d/a\), for \(a\ne0\).
\item
  \textbf{Completing the square (\texttt{completing\_square}):}
  \(x^2+bx+c=(x+b/2)^2-(b^2/4-c)\).
\item
  \textbf{Difference of squares (\texttt{difference\_of\_squares}):}
  \(a^2-b^2=(a-b)(a+b)\).
\item
  \textbf{Square of a sum (\texttt{perfect\_square\_trinomial}):}
  \((a+b)^2=a^2+2ab+b^2\); similarly \((a-b)^2=a^2-2ab+b^2\).
\item
  \textbf{Cube of a sum (\texttt{cube\_of\_binomial}):}
  \((a+b)^3=a^3+3a^2b+3ab^2+b^3\); for a difference, the two middle
  terms have negative signs.
\item
  \textbf{Sum of cubes (\texttt{sum\_of\_cubes}):}
  \(a^3+b^3=(a+b)(a^2-ab+b^2)\).
\item
  \textbf{Difference of cubes:} \(a^3-b^3=(a-b)(a^2+ab+b^2)\).
\item
  \textbf{Difference of powers (\texttt{power\_difference}):} for
  integer \(n\ge1\), \(a^n-b^n=(a-b)\sum_{k=0}^{n-1}a^{n-1-k}b^k\).
\item
  \textbf{Exponent laws (\texttt{expo\_laws}):} for integers
  \(m,n\ge0\), \(a^ma^n=a^{m+n}\) and \((a^m)^n=a^{mn}\); negative
  exponents require \(a\ne0\).
\item
  \textbf{Binomial theorem (\texttt{binomial\_expand}):} for integer
  \(n\ge0\), \((x+y)^n=\sum_{k=0}^{n}\binom nk x^{n-k}y^k\), where
  \(\binom nk=n!/(k!(n-k)!)\).
\item
  \textbf{Arithmetic progression term and sum
  (\texttt{arithmetic\_series}):} \(a_k=a+(k-1)d\) and
  \(\sum_{k=1}^{n}a_k=n(2a+(n-1)d)/2\), for integer \(n\ge0\).
\item
  \textbf{Geometric progression sum (\texttt{geometric\_series}):}
  \(\sum_{k=0}^{n-1}ar^k=a(r^n-1)/(r-1)\) for \(r\ne1\); if \(r=1\), the
  sum is \(an\).
\item
  \textbf{Gauss sum / triangular numbers:} \(\sum_{k=1}^{n}k=n(n+1)/2\),
  for integer \(n\ge0\).
\item
  \textbf{Sum of squares (\texttt{sum\_of\_squares}):}
  \(\sum_{k=1}^{n}k^2=n(n+1)(2n+1)/6\), for integer \(n\ge0\).
\item
  \textbf{Sum of cubes:} \(\sum_{k=1}^{n}k^3=[n(n+1)/2]^2\), for integer
  \(n\ge0\).
\item
  \textbf{Sum of fourth powers (\texttt{sum\_fourth\_powers}):}
  \(\sum_{k=1}^{n}k^4=n(n+1)(2n+1)(3n^2+3n-1)/30\), for integer
  \(n\ge0\).
\item
  \textbf{Sum of the first odd numbers (\texttt{sum\_first\_odds}):}
  \(\sum_{k=1}^{n}(2k-1)=n^2\), for integer \(n\ge0\).
\item
  \textbf{Telescoping reciprocal sum (\texttt{telescoping\_series}):}
  \(\sum_{k=1}^{n}1/[k(k+1)]=n/(n+1)\) because
  \(1/[k(k+1)]=1/k-1/(k+1)\).
\item
  \textbf{Horner evaluation (\texttt{horner\_eval}):}
  \(a_3x^3+a_2x^2+a_1x+a_0=((a_3x+a_2)x+a_1)x+a_0\). The rewrite
  preserves value and reduces multiplications.
\item
  \textbf{Factor theorem:} for polynomial \(p\), \(p(r)=0\iff(x-r)\)
  divides \(p(x)\). Certificate: exact polynomial division with zero
  remainder.
\item
  \textbf{Rational root theorem (\texttt{rational\_root}):} for
  integer-coefficient \(p(x)=a_nx^n+\cdots+a_0\), a reduced rational
  root \(u/v\) must have \(u\mid a_0\) and \(v\mid a_n\), with \(v>0\).
\item
  \textbf{Partial fraction identity (\texttt{partial\_fractions}):} for
  \(a\ne b\) and \(x\ne-a,-b\),
  \(1/((x+a)(x+b))=[1/(b-a)](1/(x+a)-1/(x+b))\). The current workload
  checks selected values; the engine must check the symbolic identity
  and exclusions.
\item
  \textbf{Affine composition (\texttt{composition\_affine}):} if
  \(f(x)=ax+b\) and \(g(x)=cx+d\), then \((f\circ g)(x)=acx+ad+b\).
\item
  \textbf{Inverse of a linear-fractional function
  (\texttt{inv\_function\_eval}):} for \(f(x)=ax/(x+b)\) with
  \(ab\ne0\), \(f^{-1}(y)=by/(a-y)\), with \(x\ne-b\) and \(y\ne a\).
\item
  \textbf{Absolute-value linear equation (\texttt{abs\_linear}):} for
  \(a\ne0\), \(|ax+b|=c\) has two solutions \((c-b)/a\) and \((-c-b)/a\)
  if \(c>0\), one if \(c=0\), and none if \(c<0\).
\item
  \textbf{Ratio allocation (\texttt{ratio\_share}):} if \(p,q,r\ge0\)
  and \(p+q+r>0\), splitting total \(N\) in ratio \(p:q:r\) gives
  \(Np/(p+q+r),Nq/(p+q+r),Nr/(p+q+r)\); their sum is \(N\).
\item
  \textbf{Weighted mean bounds (\texttt{weighted\_average}):} if
  \(w_i\ge0\) and \(\sum_iw_i>0\), then
  \(\min_i v_i\le\sum_iw_iv_i/\sum_iw_i\le\max_i v_i\).
\item
  \textbf{AM-GM for two nonnegative values:} \((x+y)/2\ge\sqrt{xy}\) for
  \(x,y\ge0\), since \((\sqrt{x}-\sqrt{y})^2\ge0\). Equality holds
  exactly when \(x=y\).
\item
  \textbf{Cauchy-Schwarz:} for real vectors \(u,v\),
  \((u\cdot v)^2\le(u\cdot u)(v\cdot v)\). The checker must establish
  the nonnegative difference, not just compare decimals.
\item
  \textbf{Two-square product identity (\texttt{cauchy\_identity}):}
  \((a^2+b^2)(c^2+d^2)=(ac+bd)^2+(ad-bc)^2\).
\end{itemize}

\subsection{A.3 Number-theory equations and
tests}\label{a.3-number-theory-equations-and-tests}

\begin{itemize}
\tightlist
\item
  \textbf{Least common multiple (\texttt{lcm\_pair}):} for positive
  integers \(a,b\), \(\operatorname{lcm}(a,b)=ab/\gcd(a,b)\).
\item
  \textbf{Modular exponentiation (\texttt{mod\_pow\_fast}):} maintain
  \(r b^e\equiv a^{e_0}\pmod m\); for even \(e\), square \(b\) and halve
  \(e\); for odd \(e\), multiply \(r\) by \(b\) and decrement \(e\).
  Require \(m\ge2\) and \(e\ge0\).
\item
  \textbf{Fermat's little theorem (\texttt{fermat\_little}):} if \(p\)
  is prime and \(p\nmid a\), then \(a^{p-1}\equiv1\pmod p\). This is not
  a primality test by itself.
\item
  \textbf{Euler's theorem:} if \(\gcd(a,n)=1\), then
  \(a^{\varphi(n)}\equiv1\pmod n\), where
  \(\varphi(n)=|\{1\le k\le n:\gcd(k,n)=1\}|\).
\item
  \textbf{Totient product formula:} if \(n=\prod_i p_i^{e_i}\), then
  \(\varphi(n)=n\prod_{p\mid n}(1-1/p)\), over the distinct prime
  divisors of \(n\).
\item
  \textbf{Wilson's theorem (\texttt{wilson\_check}):} integer \(p>1\) is
  prime exactly when \((p-1)!\equiv-1\pmod p\). Direct factorial
  calculation is not efficient for large \(p\).
\item
  \textbf{Chinese remainder theorem (\texttt{crt\_verify}):} for
  pairwise coprime positive moduli \(m_i\), each residue tuple has one
  solution class modulo \(M=\prod_i m_i\). Checking residues verifies a
  candidate; uniqueness also requires proving pairwise coprimality.
\item
  \textbf{Modular inverse:} \(ax\equiv1\pmod m\) has a solution iff
  \(\gcd(a,m)=1\); extended-Euclid coefficients produce one.
\item
  \textbf{Binomial coefficient recurrence (\texttt{pascal\_row\_mult}):}
  \(\binom nk=\binom n{k-1}(n-k+1)/k\) for \(1\le k\le n\); division is
  exact in integers.
\item
  \textbf{Catalan numbers:} \(C_n=\frac1{n+1}\binom{2n}{n}\) and
  \(C_n=C_{n-1}\,2(2n-1)/(n+1)\) for \(n\ge1\), with \(C_0=1\).
\item
  \textbf{Pythagorean triple identity (\texttt{pyth\_triple\_verify}):}
  for integers \(m>n>0\), \((m^2-n^2)^2+(2mn)^2=(m^2+n^2)^2\).
  Primitivity additionally requires \(\gcd(m,n)=1\) and opposite parity.
\item
  \textbf{Perfect-square test (\texttt{perfect\_square\_check}):} for
  integer \(N\ge0\), \(N\) is a square iff
  \(\lfloor\sqrt N\rfloor^2=N\). Negative integers are not integer
  squares.
\item
  \textbf{Prime trial division:} integer \(N>1\) is prime iff no integer
  \(d\) with \(2\le d\le\lfloor\sqrt N\rfloor\) divides it. A shorter
  divisor list is only a screen.
\item
  \textbf{Divisor-square sum:} for distinct primes \(p,q\) and
  \(N=p^2q\), \(\sum_{d\mid N}d^2=(1+p^2+p^4)(1+q^2)\). Distinctness
  matters because otherwise the divisor list changes.
\item
  \textbf{Fibonacci recurrence:} \(F_0=0,F_1=1,F_n=F_{n-1}+F_{n-2}\) for
  \(n\ge2\). A finite run checks only a prefix.
\item
  \textbf{Tribonacci recurrence:} \(T_n=T_{n-1}+T_{n-2}+T_{n-3}\) for
  \(n\ge3\), after all three initial values are specified.
\item
  \textbf{Collatz iteration:} \(T(n)=n/2\) for even \(n\) and
  \(T(n)=3n+1\) for odd \(n\). A fixed number of steps proves only those
  steps, not that every positive integer reaches 1.
\end{itemize}

\subsection{A.4 Linear-algebra
identities}\label{a.4-linear-algebra-identities}

\begin{itemize}
\tightlist
\item
  \textbf{Two-by-two determinant (\texttt{matrix\_det\_2x2}):}
  \(\det\begin{pmatrix}a&b\\c&d\end{pmatrix}=ad-bc\).
\item
  \textbf{Two-by-two inverse:} when \(ad-bc\ne0\),
  \(A^{-1}=\frac1{ad-bc}\begin{pmatrix}d&-b\\-c&a\end{pmatrix}\); verify
  \(AA^{-1}=I=A^{-1}A\).
\item
  \textbf{Cramer's rule (\texttt{cramer\_2x2}):} for \(Ax=b\) with
  \(\det A\ne0\), \(x_i=\det(A_i)/\det(A)\). If \(\det A=0\), analyze
  rank and consistency instead of dividing.
\item
  \textbf{Gaussian row operations:} swapping rows, multiplying a row by
  nonzero \(c\), or adding a multiple of one row to another preserves
  the solution set of the augmented system. Track row swaps and pivots.
\item
  \textbf{Rank-nullity:} for a finite-dimensional domain \(V\),
  \(\dim V=\operatorname{rank}(A)+\operatorname{nullity}(A)\).
\item
  \textbf{Determinant product:} for square matrices of the same size,
  \(\det(AB)=\det(A)\det(B)\).
\item
  \textbf{Cyclic trace:} for compatible finite matrices,
  \(\operatorname{tr}(AB)=\operatorname{tr}(BA)\). This does not imply
  \(AB=BA\).
\item
  \textbf{Rank-one trace (\texttt{outer\_trace}):}
  \(\operatorname{tr}(uv^T)=v^Tu=u\cdot v\).
\item
  \textbf{Two-dimensional Gram identity (\texttt{gram\_det}):}
  \((u\cdot u)(v\cdot v)-(u\cdot v)^2=(u_1v_2-u_2v_1)^2\ge0\).
\item
  \textbf{Parallelogram area (\texttt{area\_parallelogram2}):}
  \(\operatorname{area}=|u_1v_2-u_2v_1|\).
\item
  \textbf{Projection onto a nonzero vector:} if \(v\ne0\),
  \(\operatorname{proj}_v(u)=\frac{u\cdot v}{v\cdot v}v\). Check that
  \(u-\operatorname{proj}_v(u)\) is orthogonal to \(v\).
\item
  \textbf{Symmetric two-by-two eigenvalues (\texttt{eigen\_sym}):} for
  \(A=\begin{pmatrix}a&b\\b&d\end{pmatrix}\),
  \(\lambda_{\pm}=(a+d\pm\sqrt{(a-d)^2+4b^2})/2\). Use an exact
  algebraic or interval type if the radicand is not a rational square.
\item
  \textbf{Bareiss fraction-free update:}
  \(a_{ij}^{(k+1)}=(a_{kk}^{(k)}a_{ij}^{(k)}-a_{ik}^{(k)}a_{kj}^{(k)})/a_{k-1,k-1}^{(k-1)}\),
  with initial previous pivot 1. Require exact division and handle zero
  pivots and singular matrices.
\end{itemize}

\subsection{A.5 Calculus theorems and numerical
formulas}\label{a.5-calculus-theorems-and-numerical-formulas}

\begin{itemize}
\tightlist
\item
  \textbf{Constant and sum rules:} \((c)'=0\) and \((f+g)'=f'+g'\),
  where the derivatives exist.
\item
  \textbf{Product rule (\texttt{product\_rule}):} \((fg)'=f'g+fg'\).
\item
  \textbf{Quotient rule (\texttt{quotient\_rule}):}
  \((f/g)'=(f'g-fg')/g^2\), where \(g\ne0\) on the domain.
\item
  \textbf{Power rule (\texttt{power\_rule}):}
  \(\frac{d}{dx}x^n=nx^{n-1}\); the real-domain statement depends on
  \(n\) and \(x\), especially for negative or non-integer \(n\).
\item
  \textbf{Chain rule (\texttt{chain\_rule}):}
  \((f\circ g)'(x)=f'(g(x))g'(x)\), when both derivatives exist.
\item
  \textbf{Second derivative of a power:}
  \(\frac{d^2}{dx^2}(ax^n)=an(n-1)x^{n-2}\), where the derivatives
  exist.
\item
  \textbf{Power antiderivative:} for \(n\ne-1\),
  \(\int ax^n\,dx=ax^{n+1}/(n+1)+C\); for \(n=-1\), use \(a\ln|x|+C\) on
  intervals not crossing zero.
\item
  \textbf{Definite power integral:} for \(n\ne-1\),
  \(\int_l^h ax^n\,dx=\frac{a}{n+1}(h^{n+1}-l^{n+1})\), if the integrand
  is defined and integrable on \([l,h]\).
\item
  \textbf{Fundamental theorem of calculus:} if \(F'=f\) on \([a,b]\)
  with the required continuity/integrability assumptions,
  \(\int_a^b f(x)\,dx=F(b)-F(a)\).
\item
  \textbf{Mean value theorem:} if \(f\) is continuous on \([a,b]\) and
  differentiable on \((a,b)\), with \(a<b\), some \(c\in(a,b)\)
  satisfies \(f'(c)=(f(b)-f(a))/(b-a)\).
\item
  \textbf{Rolle's theorem:} under the MVT conditions plus \(f(a)=f(b)\),
  some \(c\in(a,b)\) has \(f'(c)=0\).
\item
  \textbf{L'Hôpital's rule:} for a differentiable quotient in a
  punctured interval, an applicable \(0/0\) or \(\infty/\infty\) form,
  and a suitable existing limit of \(f'/g'\) with \(g'\ne0\),
  \(\lim f/g=\lim f'/g'\). This is not a general quotient rewrite.
\item
  \textbf{Taylor polynomial:}
  \(f(x+h)=\sum_{k=0}^{n}f^{(k)}(x)h^k/k!+R_{n+1}\). A certified
  approximation must also bound or identify the remainder.
\item
  \textbf{Trapezoid exactness:} the one-panel trapezoid rule integrates
  every polynomial of degree at most 1 exactly.
\item
  \textbf{Simpson exactness:} the one-panel Simpson rule integrates
  every polynomial of degree at most 3 exactly, with equal subintervals
  and the standard \(1/3\) weights. Checking one quadratic does not
  prove this theorem.
\item
  \textbf{Euler method recurrence:} \(y_{k+1}=y_k+h f(x_k,y_k)\) and
  \(x_{k+1}=x_k+h\). This defines a numerical method, not an exact
  solution in general; error bounds need regularity and step-size
  assumptions.
\item
  \textbf{Newton square-root step:} for \(N>0,x_k>0\),
  \(x_{k+1}=\tfrac12(x_k+N/x_k)\). This is numerical unless exact
  divisibility or a certified interval argument applies.
\end{itemize}

\subsection{A.6 Probability and inequality
rules}\label{a.6-probability-and-inequality-rules}

\begin{itemize}
\tightlist
\item
  \textbf{Probability normalization:} \(P(\Omega)=1\), \(P(A)\ge0\), and
  for disjoint events \(A,B\), \(P(A\cup B)=P(A)+P(B)\).
\item
  \textbf{Complement and union:} \(P(A^c)=1-P(A)\) and
  \(P(A\cup B)=P(A)+P(B)-P(A\cap B)\).
\item
  \textbf{Total probability:} for finite partition \(B_1,\ldots,B_n\),
  \(P(A)=\sum_iP(A\mid B_i)P(B_i)\); do not condition on an event of
  probability zero.
\item
  \textbf{Bayes' theorem:} when \(P(B)>0\),
  \(P(A\mid B)=P(B\mid A)P(A)/P(B)\). Store the event and conditioning
  direction explicitly.
\item
  \textbf{Finite expectation:} \(\mathbb E[X]=\sum_i x_iP(X=x_i)\) with
  probabilities nonnegative and summing to one; linearity gives
  \(\mathbb E[aX+bY]=a\mathbb E[X]+b\mathbb E[Y]\).
\item
  \textbf{Variance identity:}
  \(\operatorname{Var}(X)=\mathbb E[X^2]-(\mathbb E[X])^2\ge0\), for a
  finite distribution with finite second moment.
\item
  \textbf{Triangle inequality:} \(|x+y|\le|x|+|y|\), for real or complex
  values.
\item
  \textbf{AM-GM and Cauchy-Schwarz:} use the nonnegative/real-vector
  assumptions stated in A.2; check assumptions and prove the difference
  nonnegative.
\end{itemize}

\subsection{A.7 How the current 150 entries map to this
registry}\label{a.7-how-the-current-150-entries-map-to-this-registry}

The source set contains 20 base entries, 12 I1 extensions, 68 further I1
entries, and 50 I3 additions: 150 identifiers total. The family counts
are algebra 41, calculus 36, linear algebra 33, and number theory 40.
Appendix B lists every identifier and its source-level mathematical
description. Those descriptions locate the intended workloads; an exact
sample is not a proof of the general statement written beside it.

Audit every entry before claiming theorem coverage:

\begin{enumerate}
\def\labelenumi{\arabic{enumi}.}
\tightlist
\item
  Does the formula hold over the full stated domain, or only for the
  generator's parameter range?
\item
  Are exceptional cases represented: zero denominators, repeated roots,
  negative radicands, singular matrices, non-coprime moduli, and empty
  sums?
\item
  Does the implementation check a symbolic identity or only selected
  numerical values?
\item
  Does its certificate prove the result, or only show that one candidate
  gives the expected number?
\item
  Is the item a theorem, recurrence, numerical method, finite
  computation, or test example? Keep those labels separate.
\end{enumerate}

The first implementation tranche should be the Tier 0 arithmetic rules,
linear and quadratic solving with degenerate cases, polynomial
identities and factor checks, Euclid/Bézout/modular inverse, two-by-two
exact linear algebra, and derivative/integral rules with their domains.
Add other registry entries only behind an independently checked
certificate and a measured workload.

\section{Appendix B --- Existing corpus entries, one by
one}\label{appendix-b-existing-corpus-entries-one-by-one}

This is a source-linked index of the exact 150 algorithm identifiers.
Each description is copied or lightly condensed from the
\texttt{math\_desc} attached to that identifier in code; its source link
leads to the full implementation. Treat it as the author's intended
workload description; verify it against the builder, step function,
domain assumptions, and checker before promoting it to a general
theorem. The source distribution is 20 base + 12 I1 + 68 extended I1 +
50 I3 entries.

\subsection{B.1 Algebra (41 entries)}\label{b.1-algebra-41-entries}

\begin{itemize}
\tightlist
\item
  \href{../generation_1/I3/algos_i3.py}{\texttt{abs\_linear}} ---
  \texttt{abs(a*x\ +\ b)\ =\ c\ has\ exactly\ the\ two\ solutions\ x=(c-b)/a\ and\ x=(-c-b)/a\ for\ c\textgreater{}0,\ a!=0.}
\item
  \href{../generation_1/I3/algos_i3.py}{\texttt{am\_gm\_two}} ---
  \texttt{AM-GM\ on\ the\ pair\ (a,\ 4a):\ (a+4a)/2\ \textgreater{}=\ sqrt(4a\^{}2)\ =\ 2a,\ gap\ a\ \textgreater{}=\ 0.}
\item
  \href{../tools/prm/algos.py}{\texttt{arithmetic\_series}} ---
  \texttt{sum\_\{k=0\}\^{}\{n-1\}\ (a\ +\ k*d)\ =\ n*(2a\ +\ (n-1)d)/2.}
\item
  \href{../generation_1/I1/algos_gen1_big.py}{\texttt{binomial\_coeff}}
  --- \texttt{C(n,k)\ by\ multiplicative\ formula.}
\item
  \href{../tools/prm/algos.py}{\texttt{binomial\_expand}} ---
  \texttt{(x\ +\ y)\^{}n\ =\ sum\_k\ C(n,k)\ x\^{}(n-k)\ y\^{}k.}
\item
  \href{../generation_1/I1/algos_gen1_big.py}{\texttt{cauchy\_identity}}
  --- \texttt{(a\^{}2+b\^{}2)(c\^{}2+d\^{}2)=(ac+bd)\^{}2+(ad-bc)\^{}2.}
\item
  \href{../generation_1/I1/algos_gen1_big.py}{\texttt{completing\_square}}
  --- \texttt{x\^{}2+bx+c\ =\ (x+b/2)\^{}2-((b/2)\^{}2-c).}
\item
  \href{../generation_1/I3/algos_i3.py}{\texttt{composition\_affine}}
  ---
  \texttt{f(x)=ax+b\ composed\ with\ g(x)=cx+d\ is\ (ac)x\ +\ (ad+b).}
\item
  \href{../generation_1/I1/algos_gen1_ext.py}{\texttt{cube\_of\_binomial}}
  --- \texttt{(a+b)\^{}3\ =\ a\^{}3+3a\^{}2b+3ab\^{}2+b\^{}3.}
\item
  \href{../generation_1/I1/algos_gen1_ext.py}{\texttt{difference\_of\_cubes}}
  --- \texttt{a\^{}3\ -\ b\^{}3\ =\ (a-b)(a\^{}2+ab+b\^{}2).}
\item
  \href{../tools/prm/algos.py}{\texttt{difference\_of\_squares}} ---
  \texttt{a\^{}2\ -\ b\^{}2\ =\ (a\ -\ b)(a\ +\ b).}
\item
  \href{../generation_1/I3/algos_i3.py}{\texttt{digit\_sum\_base}} ---
  \texttt{the\ base-b\ digit\ sum\ of\ N\ via\ the\ mod/div\ chain,\ exact\ for\ any\ base\ 2..8.}
\item
  \href{../generation_1/I3/algos_i3.py}{\texttt{expo\_laws}} ---
  \texttt{a\^{}m\ *\ a\^{}n\ =\ a\^{}(m+n);\ verified\ exactly\ on\ small\ integer\ powers.}
\item
  \href{../generation_1/I3/algos_i3.py}{\texttt{gauss\_pairing}} ---
  \texttt{Gauss\ pairing:\ sum\_\{1..n\}\ =\ (n/2)(n+1);\ both\ factor\ forms\ agree\ exactly.}
\item
  \href{../generation_1/I3/algos_i3.py}{\texttt{geometric\_mean\_pair}}
  ---
  \texttt{GM\ of\ (sQ\^{}2,\ s)\ equals\ sQ\ exactly;\ AM\ \textgreater{}=\ GM\ with\ gap\ (s(Q-1)\^{}2\ ...\ )/2\ \textgreater{}=\ 0.}
\item
  \href{../tools/prm/algos.py}{\texttt{geometric\_series}} ---
  \texttt{sum\_\{k=0\}\^{}\{n-1\}\ a*r\^{}k\ =\ a*(r\^{}n\ -\ 1)/(r\ -\ 1),\ for\ r\ !=\ 1.}
\item
  \href{../generation_1/I1/algos_gen1_ext.py}{\texttt{horner\_eval}} ---
  \texttt{p(x)\ =\ ((a3*x+a2)*x+a1)*x+a0.}
\item
  \href{../generation_1/I3/algos_i3.py}{\texttt{inv\_function\_eval}}
  ---
  \texttt{f(x)=ax/(x+b)\ inverts\ to\ x\ =\ yb/(a-y),\ and\ the\ round\ trip\ is\ exact.}
\item
  \href{../generation_1/I1/algos_gen1_big.py}{\texttt{linear\_recurrence}}
  --- \texttt{u\_n\ =\ p*u\_\{n-1\}+q*u\_\{n-2\},\ three\ terms.}
\item
  \href{../tools/prm/algos.py}{\texttt{linear\_solve}} ---
  \texttt{a*x\ +\ b\ =\ c\ \ =\textgreater{}\ \ x\ =\ (c\ -\ b)/a,\ for\ a\ !=\ 0.}
\item
  \href{../generation_1/I3/algos_i3.py}{\texttt{log2\_pow2}} ---
  \texttt{if\ N\ is\ a\ power\ of\ two,\ N\ =\ 2\^{}k\ and\ k\ is\ the\ halving\ count.}
\item
  \href{../generation_1/I1/algos_gen1_ext.py}{\texttt{log\_product}} ---
  \texttt{log2(a*b)\ =\ log2(a)+log2(b),\ a,b\ powers\ of\ two.}
\item
  \href{../generation_1/I3/algos_i3.py}{\texttt{median\_five}} ---
  \texttt{median\ of\ five\ distinct\ numbers:\ two\ pair\ partitions\ leave\ three\ candidates;\ sum-min-max\ is\ the\ middle.}
\item
  \href{../generation_1/I1/algos_gen1_big.py}{\texttt{nested\_radical}}
  --- \texttt{(sqrt(a)+sqrt(b))\^{}2\ =\ a+b+2sqrt(ab).}
\item
  \href{../generation_1/I1/algos_gen1_big.py}{\texttt{partial\_fractions}}
  --- \texttt{1/((x+a)(x+b))\ split\ verified\ at\ x0.}
\item
  \href{../generation_1/I3/algos_i3.py}{\texttt{percentage\_pair}} ---
  \texttt{q\%\ of\ n\ is\ qn/100\ and\ reversing\ recovers\ n\ exactly.}
\item
  \href{../tools/prm/algos.py}{\texttt{perfect\_square\_trinomial}} ---
  \texttt{a\^{}2\ +\ 2ab\ +\ b\^{}2\ =\ (a\ +\ b)\^{}2.}
\item
  \href{../generation_1/I1/algos_gen1_big.py}{\texttt{power\_difference}}
  --- \texttt{a\^{}n-b\^{}n\ =\ (a-b)*sum,\ n\textless{}=4.}
\item
  \href{../tools/prm/algos.py}{\texttt{quadratic\_formula}} ---
  \texttt{a*x\^{}2\ +\ b*x\ +\ c\ =\ 0\ \ =\textgreater{}\ \ x\ =\ (-b\ +-\ sqrt(b\^{}2-4ac))\ /\ (2a).}
\item
  \href{../generation_1/I1/algos_gen1_big.py}{\texttt{quartic\_horner}}
  --- \texttt{Degree-4\ Horner\ in\ 3\ stages.}
\item
  \href{../generation_1/I3/algos_i3.py}{\texttt{ratio\_share}} ---
  \texttt{N\ split\ in\ ratio\ p:q:r\ gives\ unit\ N/(p+q+r)\ and\ shares\ unit*p,\ unit*q,\ unit*r.}
\item
  \href{../generation_1/I1/algos_gen1_big.py}{\texttt{rational\_root}}
  --- \texttt{p(r)=0\ via\ Horner\ for\ (x-r)(ax+b).}
\item
  \href{../tools/prm/algos.py}{\texttt{slope\_from\_points}} ---
  \texttt{Through\ (x1,y1),(x2,y2):\ m\ =\ (y2-y1)/(x2-x1),\ b\ =\ y1\ -\ m*x1,\ so\ y\ =\ m*x\ +\ b.}
\item
  \href{../tools/prm/algos.py}{\texttt{substitution\_system}} ---
  \texttt{x\ +\ y\ =\ S\ and\ x\ -\ y\ =\ D\ \ =\textgreater{}\ \ x\ =\ (S+D)/2,\ y\ =\ (S-D)/2.}
\item
  \href{../generation_1/I3/algos_i3.py}{\texttt{sum\_first\_odds}} ---
  \texttt{the\ sum\ of\ the\ first\ n\ odd\ numbers\ is\ n\^{}2:\ average\ (2n-1+1)/2\ =\ n,\ times\ n\ terms.}
\item
  \href{../generation_1/I1/algos_gen1_big.py}{\texttt{sum\_fourth\_powers}}
  --- \texttt{sum\ k\^{}4\ =\ n(n+1)(2n+1)(3n\^{}2+3n-1)/30.}
\item
  \href{../generation_1/I1/algos_gen1_big.py}{\texttt{sum\_odd}} ---
  \texttt{Sum\ of\ first\ odds\ telescopes\ in\ partials.}
\item
  \href{../tools/prm/algos.py}{\texttt{sum\_of\_cubes}} ---
  \texttt{a\^{}3\ +\ b\^{}3\ =\ (a\ +\ b)(a\^{}2\ -\ ab\ +\ b\^{}2).}
\item
  \href{../generation_1/I1/algos_gen1_ext.py}{\texttt{sum\_of\_squares}}
  --- \texttt{1\^{}2\ +\ ...\ +\ n\^{}2\ =\ n(n+1)(2n+1)/6.}
\item
  \href{../generation_1/I1/algos_gen1_ext.py}{\texttt{vieta\_cubic}} ---
  \texttt{For\ roots\ r1,r2,r3:\ sum,\ pairwise\ sum,\ product.}
\item
  \href{../generation_1/I3/algos_i3.py}{\texttt{weighted\_average}} ---
  \texttt{weighted\ mean\ (k1\ v1\ +\ k2\ v2)/(k1\ +\ k2)\ lies\ between\ v1\ and\ v2,\ computed\ exactly.}
\end{itemize}

\subsection{B.2 Calculus (36 entries)}\label{b.2-calculus-36-entries}

\begin{itemize}
\tightlist
\item
  \href{../generation_1/I3/algos_i3.py}{\texttt{area\_under\_linear}}
  ---
  \texttt{area\ under\ mx+c\ on\ {[}a,b{]}:\ m(b\^{}2-a\^{}2)/2\ +\ c(b-a),\ exact.}
\item
  \href{../generation_1/I3/algos_i3.py}{\texttt{asymptote\_oblique}} ---
  \texttt{divide\ (Ax\^{}2+Bx+C)\ by\ (x-d):\ quotient\ A\ x\ +\ (B+Ad),\ remainder\ C+d(B+Ad),\ vertical\ at\ d.}
\item
  \href{../generation_1/I1/algos_gen1_big.py}{\texttt{average\_linear}}
  --- \texttt{Average\ of\ a*x+b\ on\ {[}0,2{]}\ is\ a+b.}
\item
  \href{../generation_1/I1/algos_gen1_big.py}{\texttt{chain\_poly}} ---
  \texttt{d/dx\ (ax+b)\^{}n\ at\ x0.}
\item
  \href{../tools/prm/algos.py}{\texttt{chain\_rule}} ---
  \texttt{d/dx\ (a*x\^{}n\ +\ b)\ =\ a*n*x\^{}(n-1);\ the\ constant\ b\ differentiates\ to\ 0.}
\item
  \href{../generation_1/I1/algos_gen1_big.py}{\texttt{concavity\_quad}}
  ---
  \texttt{f\textquotesingle{}\textquotesingle{}\ =\ 2a\ sign\ for\ a*x\^{}2.}
\item
  \href{../generation_1/I3/algos_i3.py}{\texttt{critical\_cubic}} ---
  \texttt{critical\ points\ of\ x\^{}3+ax\^{}2+bx+c\ via\ the\ perfect-square\ discriminant\ of\ f\textquotesingle{}.}
\item
  \href{../tools/prm/algos.py}{\texttt{definite\_integral}} ---
  \texttt{int\_lo\^{}hi\ a*x\^{}n\ dx\ =\ a/(n+1)\ *\ (hi\^{}(n+1)\ -\ lo\^{}(n+1)).}
\item
  \href{../generation_1/I1/algos_gen1_big.py}{\texttt{definite\_quad}}
  --- \texttt{int\_0\^{}1\ (ax\^{}2+bx+c).}
\item
  \href{../generation_1/I3/algos_i3.py}{\texttt{euler\_method\_linear}}
  ---
  \texttt{Euler\ method\ for\ y\textquotesingle{}=kx,\ y(x0)=y0:\ two\ steps\ with\ rational\ h\ are\ exact.}
\item
  \href{../generation_1/I1/algos_gen1_big.py}{\texttt{ftc\_check}} ---
  \texttt{F(1)-F(0)\ for\ a*x\^{}n.}
\item
  \href{../generation_1/I1/algos_gen1_ext.py}{\texttt{harmonic\_partial}}
  ---
  \texttt{H\_n\ =\ sum\_\{k=1..n\}\ 1/k\ exact,\ n\ \textless{}=\ 6.}
\item
  \href{../generation_1/I1/algos_gen1_big.py}{\texttt{implicit\_circle}}
  --- \texttt{Circle\ slope\ at\ (0,r)\ is\ 0.}
\item
  \href{../generation_1/I1/algos_gen1_big.py}{\texttt{inflection\_cubic}}
  --- \texttt{x\^{}3+b*x\ inflects\ at\ 0.}
\item
  \href{../generation_1/I1/algos_gen1_big.py}{\texttt{inverse\_slope}}
  --- \texttt{Inverse\ slope\ 1/m.}
\item
  \href{../generation_1/I1/algos_gen1_big.py}{\texttt{lhopital\_poly}}
  --- \texttt{lim\ (x\^{}2-a\^{}2)/(x-a)\ =\ 2a.}
\item
  \href{../generation_1/I1/algos_gen1_big.py}{\texttt{monotonic\_quad}}
  --- \texttt{Derivative\ of\ a*x\^{}2+x\ at\ 0\ and\ 1.}
\item
  \href{../generation_1/I3/algos_i3.py}{\texttt{mvt\_quadratic}} ---
  \texttt{Mean\ value\ theorem\ for\ a\ quadratic:\ the\ tangent-slope\ point\ is\ exactly\ (p+q)/2.}
\item
  \href{../generation_1/I3/algos_i3.py}{\texttt{newton\_sqrt\_step}} ---
  \texttt{one\ Newton\ step\ for\ sqrt(N):\ x1=(x0+N/x0)/2,\ kept\ exact\ because\ x0\ \textbar{}\ N.}
\item
  \href{../tools/prm/algos.py}{\texttt{power\_rule}} ---
  \texttt{d/dx\ (a*x\^{}n)\ =\ a*n*x\^{}(n-1).}
\item
  \href{../generation_1/I3/algos_i3.py}{\texttt{power\_rule\_integrate}}
  ---
  \texttt{integral\ of\ c\ x\^{}n\ over\ {[}0,B{]}\ is\ c\ B\^{}(n+1)/(n+1),\ exact\ in\ rationals.}
\item
  \href{../tools/prm/algos.py}{\texttt{product\_rule}} ---
  \texttt{d/dx\ {[}\ (a\ x\^{}m)(b\ x\^{}n)\ {]}\ =\ a*m*b*x\^{}(n+m-1)\ +\ a*b*n*x\^{}(m+n-1).}
\item
  \href{../tools/prm/algos.py}{\texttt{quotient\_rule}} ---
  \texttt{d/dx\ {[}\ u/v\ {]}\ =\ (u\textquotesingle{}v\ -\ uv\textquotesingle{})/v\^{}2,\ with\ u\ =\ a\ x\^{}2\ +\ b,\ v\ =\ c\ x\ +\ d\ evaluated\ at\ x\ =\ 1.}
\item
  \href{../generation_1/I1/algos_gen1_big.py}{\texttt{riemann\_mid2}}
  --- \texttt{2-rectangle\ midpoint\ of\ x\^{}2\ on\ {[}0,2{]}.}
\item
  \href{../generation_1/I3/algos_i3.py}{\texttt{rolle\_check}} ---
  \texttt{f=A(x-p)(x-q):\ f\ vanishes\ at\ p,q\ and\ f\textquotesingle{}\ vanishes\ exactly\ at\ (p+q)/2,\ strictly\ between.}
\item
  \href{../generation_1/I3/algos_i3.py}{\texttt{secant\_poly\_slope}}
  ---
  \texttt{secant\ slope\ of\ Ax\^{}2+B\ between\ x1,\ x2\ equals\ A(x1+x2)\ exactly.}
\item
  \href{../generation_1/I1/algos_gen1_big.py}{\texttt{second\_derivative}}
  --- \texttt{Second\ derivative\ of\ a*x\^{}n\ at\ x0.}
\item
  \href{../generation_1/I3/algos_i3.py}{\texttt{second\_diff\_central}}
  ---
  \texttt{central\ second\ difference\ of\ Ax\^{}3\ equals\ 6Ax\ exactly,\ the\ h\^{}2\ error\ term\ vanishing\ for\ cubics.}
\item
  \href{../generation_1/I1/algos_gen1_big.py}{\texttt{shadow\_rates}}
  --- \texttt{Similar-triangle\ shadow\ =\ 2h.}
\item
  \href{../generation_1/I3/algos_i3.py}{\texttt{simpson\_quadratic}} ---
  \texttt{Simpson\textquotesingle{}s\ rule\ is\ exact\ for\ quadratics:\ S\ -\ I\ =\ 0,\ verified\ exactly.}
\item
  \href{../generation_1/I3/algos_i3.py}{\texttt{tangent\_line\_quad}}
  ---
  \texttt{tangent\ to\ ax\^{}2+bx+c\ at\ x0:\ slope\ 2ax0+b,\ intercept\ c\ -\ a\ x0\^{}2.}
\item
  \href{../generation_1/I1/algos_gen1_big.py}{\texttt{tangent\_point}}
  --- \texttt{Value\ and\ slope\ of\ a*x\^{}n\ at\ x0.}
\item
  \href{../generation_1/I1/algos_gen1_big.py}{\texttt{taylor\_linear}}
  --- \texttt{Linear\ approx\ of\ a*x\^{}2\ at\ x0+h.}
\item
  \href{../generation_1/I1/algos_gen1_ext.py}{\texttt{telescoping\_series}}
  --- \texttt{sum\_\{k=1..n\}\ 1/(k(k+1))\ =\ n/(n+1).}
\item
  \href{../generation_1/I1/algos_gen1_big.py}{\texttt{trapezoid\_linear}}
  --- \texttt{Trapezoid\ of\ a*x+b\ on\ {[}0,2{]}\ is\ exact.}
\item
  \href{../generation_1/I1/algos_gen1_big.py}{\texttt{usub\_linear}} ---
  \texttt{d/dx\ (ax+b)\^{}n\ at\ x=1.}
\end{itemize}

\subsection{B.3 Linear algebra (33
entries)}\label{b.3-linear-algebra-33-entries}

\begin{itemize}
\tightlist
\item
  \href{../generation_1/I3/algos_i3.py}{\texttt{affine\_point}} ---
  \texttt{r\ =\ p\ +\ alpha(q-p)\ is\ the\ affine\ interpolation\ point\ on\ segment\ pq,\ exact\ in\ rationals.}
\item
  \href{../generation_1/I3/algos_i3.py}{\texttt{area\_parallelogram2}}
  ---
  \texttt{parallelogram\ area\ from\ column\ vectors\ is\ \textbar{}a1\ b2\ -\ a2\ b1\textbar{},\ exact.}
\item
  \href{../generation_1/I1/algos_gen1_big.py}{\texttt{commutator\_trace}}
  --- \texttt{tr(AB)-tr(BA)\ =\ 0.}
\item
  \href{../generation_1/I1/algos_gen1_big.py}{\texttt{cramer\_2x2}} ---
  \texttt{Cramer\ solution\ of\ 2x2.}
\item
  \href{../generation_1/I1/algos_gen1_big.py}{\texttt{cross\_2d}} ---
  \texttt{2-D\ cross\ a1b2-a2b1.}
\item
  \href{../generation_1/I3/algos_i3.py}{\texttt{dep\_pair\_cross}} ---
  \texttt{two\ vectors\ are\ parallel\ exactly\ when\ their\ 2D\ cross\ product\ is\ zero;\ the\ ratio\ is\ the\ scale.}
\item
  \href{../generation_1/I3/algos_i3.py}{\texttt{det\_3x3\_cofactor}} ---
  \texttt{det\ of\ a\ 3x3\ by\ the\ six\ Sarrus\ triple\ products,\ exact\ integers.}
\item
  \href{../generation_1/I1/algos_gen1_big.py}{\texttt{det\_product}} ---
  \texttt{det(AB)\ from\ both\ sides.}
\item
  \href{../generation_1/I1/algos_gen1_big.py}{\texttt{diag\_trace}} ---
  \texttt{Trace\ of\ a\ squared\ diagonal.}
\item
  \href{../generation_1/I3/algos_i3.py}{\texttt{diagonal\_eigencheck}}
  ---
  \texttt{a\ diagonal\ matrix\ has\ its\ eigenvalues\ exactly\ on\ the\ diagonal;\ det\ and\ trace\ confirm.}
\item
  \href{../tools/prm/algos.py}{\texttt{dot\_product\_3d}} ---
  \texttt{a\ .\ b\ =\ a1*b1\ +\ a2*b2\ +\ a3*b3.}
\item
  \href{../generation_1/I1/algos_gen1_big.py}{\texttt{eigen\_sym}} ---
  \texttt{Eigenvalues\ of\ symmetric\ 2x2,\ exact.}
\item
  \href{../tools/prm/algos.py}{\texttt{gauss\_2x2}} ---
  \texttt{Solve\ {[}{[}a,b{]},{[}c,d{]}{]}\ x\ =\ {[}e,f{]}:\ det\ =\ ad-bc,\ x1\ =\ (de-bf)/det,\ x2\ =\ (af-ce)/det.}
\item
  \href{../generation_1/I1/algos_gen1_big.py}{\texttt{gram\_det}} ---
  \texttt{Gram\ determinant\ identity.}
\item
  \href{../generation_1/I1/algos_gen1_big.py}{\texttt{hadamard\_trace}}
  --- \texttt{Entrywise-square\ trace.}
\item
  \href{../generation_1/I1/algos_gen1_big.py}{\texttt{idempotent\_tr}}
  --- \texttt{Rank-1\ projector\ trace\ is\ 1.}
\item
  \href{../generation_1/I1/algos_gen1_big.py}{\texttt{inverse\_2x2}} ---
  \texttt{det\ and\ diagonal\ of\ the\ inverse.}
\item
  \href{../generation_1/I3/algos_i3.py}{\texttt{inverse\_mult\_identity2}}
  ---
  \texttt{A\^{}\{-1\}\ times\ A\ is\ the\ identity:\ all\ four\ product\ entries\ are\ det/det,\ 0,\ 0,\ det/det.}
\item
  \href{../generation_1/I1/algos_gen1_big.py}{\texttt{mat\_add}} ---
  \texttt{Entrywise\ sums\ s11,\ s22.}
\item
  \href{../generation_1/I1/algos_gen1_big.py}{\texttt{mat\_square}} ---
  \texttt{Diagonal\ of\ A\^{}2.}
\item
  \href{../tools/prm/algos.py}{\texttt{matrix\_det\_2x2}} ---
  \texttt{det({[}{[}a,b{]},{[}c,d{]}{]})\ =\ ad\ -\ bc.}
\item
  \href{../generation_1/I3/algos_i3.py}{\texttt{matrix\_sym\_part}} ---
  \texttt{the\ symmetric\ part\ (A+A\^{}T)/2\ of\ a\ 2x2\ averages\ each\ off-diagonal\ pair\ exactly.}
\item
  \href{../generation_1/I1/algos_gen1_ext.py}{\texttt{matrix\_trace\_product}}
  --- \texttt{tr(AB)\ =\ (AB)\_11\ +\ (AB)\_22\ for\ 2x2.}
\item
  \href{../generation_1/I3/algos_i3.py}{\texttt{matvec\_3}} ---
  \texttt{3x3\ matrix\ times\ a\ 3-vector,\ three\ exact\ dot\ products.}
\item
  \href{../generation_1/I1/algos_gen1_big.py}{\texttt{nilpotent\_sq}}
  --- \texttt{Strictly-upper\ 2x2\ squares\ to\ zero.}
\item
  \href{../generation_1/I1/algos_gen1_big.py}{\texttt{outer\_trace}} ---
  \texttt{tr(uvT)\ =\ u.v.}
\item
  \href{../generation_1/I3/algos_i3.py}{\texttt{proj\_pyth}} ---
  \texttt{orthogonal\ projection\ of\ u\ onto\ a\ Pythagorean\ vector\ (m,n,\textbar{}v\textbar{}=s)\ with\ exact\ coef\ uv/\textbar{}v\textbar{}\^{}2.}
\item
  \href{../generation_1/I1/algos_gen1_big.py}{\texttt{project\_axis}}
  --- \texttt{Projection\ onto\ a\ coordinate\ axis.}
\item
  \href{../generation_1/I1/algos_gen1_big.py}{\texttt{rotate\_twice}}
  --- \texttt{Two\ 90-degree\ rotations\ negate.}
\item
  \href{../generation_1/I3/algos_i3.py}{\texttt{row\_reduce\_rank2}} ---
  \texttt{rank\ of\ a\ 2x3:\ one\ exact\ elimination\ row\ op,\ then\ count\ the\ free\ columns.}
\item
  \href{../generation_1/I1/algos_gen1_big.py}{\texttt{scalar\_mult}} ---
  \texttt{Scalar\ multiples.}
\item
  \href{../generation_1/I3/algos_i3.py}{\texttt{trace\_frobenius2}} ---
  \texttt{Frobenius\ norm\ squared\ of\ a\ 2x2:\ four\ entry\ squares\ summed,\ exact.}
\item
  \href{../generation_1/I3/algos_i3.py}{\texttt{vector\_angle\_cossq}}
  ---
  \texttt{cos\^{}2\ of\ the\ angle\ between\ Pythagorean\ vectors\ is\ (u.v)\^{}2/(\textbar{}u\textbar{}\^{}2\textbar{}v\textbar{}\^{}2),\ fully\ rational.}
\end{itemize}

\subsection{B.4 Number theory (40
entries)}\label{b.4-number-theory-40-entries}

\begin{itemize}
\tightlist
\item
  \href{../generation_1/I3/algos_i3.py}{\texttt{bayes\_two\_events}} ---
  \texttt{two-way\ contingency\ table\ with\ exact\ prior\ and\ posterior\ probabilities.}
\item
  \href{../generation_1/I1/algos_gen1_big.py}{\texttt{bezout\_witness}}
  --- \texttt{Small\ Bezout\ witnesses.}
\item
  \href{../generation_1/I1/algos_gen1_big.py}{\texttt{binary\_value}}
  --- \texttt{Binary\ digit\ accumulation.}
\item
  \href{../generation_1/I3/algos_i3.py}{\texttt{catalan\_small}} ---
  \texttt{Catalan\ numbers\ by\ the\ exact\ multiplicative\ recurrence\ C\_k\ =\ C\_\{k-1\}\ *\ 2(2k-1)/(k+1).}
\item
  \href{../generation_1/I1/algos_gen1_big.py}{\texttt{collatz\_5}} ---
  \texttt{First\ five\ Collatz\ iterates.}
\item
  \href{../generation_1/I1/algos_gen1_big.py}{\texttt{contfrac\_2}} ---
  \texttt{a0+1/(a1+1/a2)\ exact.}
\item
  \href{../generation_1/I1/algos_gen1_big.py}{\texttt{crt\_verify}} ---
  \texttt{Residues\ of\ a\ CRT\ solution.}
\item
  \href{../generation_1/I3/algos_i3.py}{\texttt{digit\_flipper}} ---
  \texttt{reverse\ a\ three-digit\ number\ through\ its\ mod/div\ digits;\ then\ take\ the\ signed\ difference.}
\item
  \href{../generation_1/I1/algos_gen1_big.py}{\texttt{digital\_root}}
  --- \texttt{Digit\ sum\ and\ mod-9\ root.}
\item
  \href{../generation_1/I1/algos_gen1_big.py}{\texttt{divisor\_count}}
  --- \texttt{tau(n)\ by\ counting.}
\item
  \href{../generation_1/I3/algos_i3.py}{\texttt{divisor\_sq\_two\_primes}}
  ---
  \texttt{sigma\_2\ of\ N=p\^{}2\ q\ sums\ d\^{}2\ over\ all\ six\ divisors,\ exact\ while\ p,q\ stay\ prime.}
\item
  \href{../tools/prm/algos.py}{\texttt{euclid\_gcd}} ---
  \texttt{gcd(x,y)\ by\ repeated\ remainder:\ gcd(x,y)\ =\ gcd(y,\ x\ mod\ y)\ until\ the\ remainder\ is\ 0.}
\item
  \href{../generation_1/I1/algos_gen1_big.py}{\texttt{euler\_phi}} ---
  \texttt{phi(n)\ by\ counting,\ small\ n.}
\item
  \href{../generation_1/I1/algos_gen1_ext.py}{\texttt{factorial\_recurrence}}
  --- \texttt{n!\ by\ recurrence,\ n\ \textless{}=\ 7.}
\item
  \href{../generation_1/I1/algos_gen1_big.py}{\texttt{fermat\_little}}
  --- \texttt{a\^{}(p-1)\ =\ 1\ mod\ p.}
\item
  \href{../generation_1/I1/algos_gen1_ext.py}{\texttt{fibonacci\_iter}}
  --- \texttt{Fibonacci\ by\ iteration\ to\ F\_n,\ n\ \textless{}=\ 10.}
\item
  \href{../generation_1/I1/algos_gen1_big.py}{\texttt{gcd\_three}} ---
  \texttt{gcd(a,b,c)\ in\ two\ steps.}
\item
  \href{../generation_1/I1/algos_gen1_big.py}{\texttt{lcm\_pair}} ---
  \texttt{lcm\ =\ a*b/gcd.}
\item
  \href{../generation_1/I1/algos_gen1_big.py}{\texttt{legendre\_euler}}
  --- \texttt{Euler\ criterion\ power\ mod\ p.}
\item
  \href{../tools/prm/algos.py}{\texttt{mod\_pow\_fast}} ---
  \texttt{base\^{}exp\ mod\ m\ by\ binary\ (square-and-multiply)\ exponentiation.}
\item
  \href{../generation_1/I1/algos_gen1_ext.py}{\texttt{modular\_inverse}}
  --- \texttt{a*x\ =\ 1\ (mod\ m)\ for\ coprime\ a,m.}
\item
  \href{../generation_1/I1/algos_gen1_big.py}{\texttt{mult\_order}} ---
  \texttt{Multiplicative\ order\ mod\ m.}
\item
  \href{../generation_1/I1/algos_gen1_big.py}{\texttt{multiple\_six}}
  --- \texttt{6k\ mod\ 6\ is\ 0.}
\item
  \href{../generation_1/I3/algos_i3.py}{\texttt{pascal\_row\_mult}} ---
  \texttt{Pascal\ row\ n\ by\ the\ exact\ multiplicative\ formula\ C(n,k)=C(n,k-1)(n-k+1)/k.}
\item
  \href{../generation_1/I1/algos_gen1_big.py}{\texttt{pell\_check}} ---
  \texttt{x\^{}2-Dy\^{}2=1\ for\ tabulated\ minimal\ solutions.}
\item
  \href{../generation_1/I3/algos_i3.py}{\texttt{perfect\_square\_check}}
  ---
  \texttt{N\ is\ a\ perfect\ square\ exactly\ when\ isqrt(N)\^{}2\ ==\ N;\ verified\ in\ integers.}
\item
  \href{../generation_1/I3/algos_i3.py}{\texttt{popcount\_halve}} ---
  \texttt{popcount\ of\ N\textless{}=31\ by\ five\ halving\ rounds,\ exact\ bit\ arithmetic.}
\item
  \href{../generation_1/I1/algos_gen1_big.py}{\texttt{power\_two}} ---
  \texttt{2\^{}e\ and\ its\ parity.}
\item
  \href{../generation_1/I1/algos_gen1_big.py}{\texttt{prime\_divisor}}
  --- \texttt{A\ divisor\ and\ zero\ remainder.}
\item
  \href{../generation_1/I3/algos_i3.py}{\texttt{prime\_factor\_small}}
  ---
  \texttt{factor\ N=p\^{}a\ q\^{}b\ by\ trial\ division\ through\ the\ small\ primes,\ exact.}
\item
  \href{../generation_1/I3/algos_i3.py}{\texttt{prob\_urn}} ---
  \texttt{drawing\ from\ an\ urn\ with\ r\ red\ and\ b\ blue:\ P\ values\ are\ exact\ rationals\ summing\ to\ 1.}
\item
  \href{../generation_1/I3/algos_i3.py}{\texttt{pyth\_triple\_verify}}
  ---
  \texttt{Euclid\textquotesingle{}s\ formula:\ legs\ 2mn\ and\ m\^{}2-n\^{}2\ give\ hypotenuse\ m\^{}2+n\^{}2,\ verified\ to\ c\^{}2=a\^{}2+b\^{}2.}
\item
  \href{../generation_1/I1/algos_gen1_big.py}{\texttt{repunit\_mod}} ---
  \texttt{Repunit\ residue.}
\item
  \href{../generation_1/I1/algos_gen1_big.py}{\texttt{sigma\_small}} ---
  \texttt{Divisor-sum\ partials.}
\item
  \href{../generation_1/I1/algos_gen1_big.py}{\texttt{square\_part}} ---
  \texttt{n\ =\ d*s\^{}2\ decomposition.}
\item
  \href{../generation_1/I1/algos_gen1_big.py}{\texttt{triangular\_mod}}
  --- \texttt{Triangular\ number\ and\ 8T+1.}
\item
  \href{../generation_1/I1/algos_gen1_big.py}{\texttt{triangular\_verify}}
  --- \texttt{T\_n\ by\ recurrence\ and\ closed\ form.}
\item
  \href{../generation_1/I3/algos_i3.py}{\texttt{tribonacci\_general}}
  ---
  \texttt{generalized\ tribonacci\ T(n)=T(n-1)+T(n-2)+T(n-3)\ from\ starts\ a,b,c,\ six\ exact\ terms.}
\item
  \href{../generation_1/I3/algos_i3.py}{\texttt{twin\_prime\_screen}}
  ---
  \texttt{twin-prime\ screen:\ p,\ p+2\ checked\ against\ small-divisor\ residues,\ verdict\ from\ primality.}
\item
  \href{../generation_1/I1/algos_gen1_big.py}{\texttt{wilson\_check}}
  --- \texttt{(p-1)!\ =\ p-1\ mod\ p.}
\end{itemize}

\end{document}
